% Point où les champs de gravitation de la Terre et de la Lune se compensent (corrigé de l'exercice 11, chapitre 11)
\begin{tikzpicture}[x=1cm,y=1cm]
\shade[ball color=otNavy!35] (0,0) circle (0.7); \node[pcLegendeG] at (0,0) {T};
\shade[ball color=otMuted!40] (10,0) circle (0.25); \node[pcLegende,anchor=north] at (10,-0.3) {L};
\draw[otMuted,line width=0.4pt] (0.7,0) -- (9.75,0);
\fill[otInk] (9.0,0) circle (0.06); \node[pcLegendeG,anchor=south] at (9.0,0.05) {P};
\draw[otNavy,line width=1.2pt,-{Stealth[length=2.4mm]}] (9.0,0.0) -- (7.9,0.0) node[pcLegende,anchor=south,text=otNavy] {$\vec{g}_\mathrm{T}$};
\draw[otGreenDark,line width=1.2pt,-{Stealth[length=2.4mm]}] (9.0,-0.02) -- (10.1,-0.02);
\node[pcLegende,anchor=north,text=otGreenDark] at (9.55,-0.1) {$\vec{g}_\mathrm{L}$};
\draw[otMuted,line width=0.4pt,{Stealth[length=1.4mm]}-{Stealth[length=1.4mm]}] (0,-1.0) -- (9.0,-1.0) node[midway,below,pcLegende] {$x\approx0{,}90\,d_\mathrm{TL}\approx\qty{3,46e8}{m}$};
\draw[otMuted,line width=0.4pt,{Stealth[length=1.4mm]}-{Stealth[length=1.4mm]}] (9.0,-1.0) -- (10,-1.0) node[midway,below,pcLegende] {$d_\mathrm{TL}-x$};
\end{tikzpicture}
\par\smallskip\legende{En P, les deux champs ont la même valeur et des sens opposés : leur somme est nulle. P est neuf fois plus loin de la Terre que de la Lune.}
